\section{Giriş}
\label{sec:intro}

Gastrointestinal (GI) endoskopi görüntülerinin otomatik sınıflandırılması,
tarama ve tanı süreçlerinde uzman iş yükünü azaltma potansiyeli taşıyan bir
bilgisayarlı görü problemidir. Bu çalışmada, halka açık \textbf{HiperKvasir}
veri kümesinin 23 sınıflı etiketli alt kümesi üzerinde, anatomik dönüm
noktaları, patolojik bulgular ve tedavi/aletsel durumları kapsayan çok sınıflı
bir sınıflandırma görevi ele alınmaktadır~\cite{borgli2020hyperkvasir}. Veri
kümesinin başlıca zorluğu, sınıf dağılımının \emph{aşırı dengesiz} olmasıdır:
bazı sınıflar binlerce görüntüye sahipken, bazı nadir sınıflar yalnızca onlu
sayılarda örnek içerir.

Önerilen yaklaşım, tek bir omurga yerine \emph{birden çok görüntü dönüştürücü
(Vision Transformer, ViT)} omurgasını özellik çıkarıcı olarak kullanıp bunların
özniteliklerini birleştiren bir özellik füzyonu mimarisidir. Üç farklı tümleyici
omurga seçilmiştir: saf yerel-olmayan dikkat temelli \textbf{ViT-B/16}, hiyerarşik
kaydırmalı-pencere temelli \textbf{Swin-T} ve maskeli görüntü modellemesiyle
öz-denetimli ön-eğitilmiş \textbf{BEiT-B/16}. Bu üçlü; küresel öz-dikkat (ViT),
hiyerarşik yerel-çok ölçekli temsil (Swin) ve öz-denetimli ön-eğitim (BEiT) gibi
farklı tümdengelimsel önyargıları bir araya getirir; böylece tek bir mimari ailesine
bağlı kalmadan tamamlayıcı temsiller öğrenmeyi amaçlar.

Çalışmanın karşılaştırma eksenleri üç düzeydedir: (i)~\emph{omurga düzeyinde}
(tekli/ikili/üçlü kombinasyonlar), (ii)~\emph{füzyon düzeyinde} (birleştirme,
ağırlıklı ve kapılı GMU füzyonu) ve (iii)~\emph{transfer öğrenme düzeyinde}
(dondurulmuş öznitelik çıkarımı ile son blokların ince ayarı). Sınıflandırıcı
tüm yapılandırmalarda sabit tutulmuştur (tek bir MLP); dolayısıyla
sınıflandırıcı-düzeyinde bir karşılaştırma \emph{yapılmamıştır}, böylece
gözlemlenen tüm farklar omurga, füzyon ve transfer seçimlerine atfedilebilir.

Ana katkı ve bulgular şu şekilde özetlenebilir:
\begin{itemize}
  \item Üç ViT omurgasının ağırlıklı füzyonu, resmi 5-katlı protokolde en iyi
        makro-F1 değerini vermiştir; bu yapılandırma eşli McNemar testinde en iyi
        tekli omurgadan ve en iyi ikili füzyondan anlamlı biçimde daha doğrudur.
  \item Gelişmiş kapılı füzyon (GMU), ağırlıklı füzyona kıyasla bir kazanım
        sağlamamıştır.
  \item Modelin başarısı, dengesizlik nedeniyle bir tavana ulaşmıştır: yüksek
        doğruluk (\%88) ile daha düşük makro-F1 ($\approx 0{,}61$) arasındaki fark,
        nadir sınıfların başarımı sınırladığını açıkça gösterir. Bu olgu, dikkat
        haritaları, Grad-CAM++ ve UMAP görselleştirmeleriyle de desteklenmiştir.
\end{itemize}
